% =========================================================================
% Results chapter: system integration and climate sensitivity (RQ3 + RQ4).
% Drop-in for the four-section skeleton of "Results: System Integration
% and Climate Sensitivity" in the msc-thesis repo. Every number traces to
% the repository campaign outputs (results/system_integration/phase3/,
% commit-pinned) produced by scripts/phase3_campaign.py under ADR-0038/
% 0042/0043; the surface-curve derivation decisions and the two model
% findings are recorded in ADR-0042 (repository decision log).
%
% Existing BibLaTeX keys reused: uemura_phd_2025, uemura_wp2_2024,
%   pol_ress_2023, dean_2010, kanning_2012, hoffmans_2014,
%   schweckendiek_2014.
% No new keys required.
% Figures (canonical copies in the engine repo under docs/figures/):
%   figures/phase3_system_fragility_bep_sections.png
%   figures/phase3_dominance_profile.png
%   figures/phase3_climate_shift.png
%   figures/phase3_rq4_attribution.png
%   figures/phase3_event_based_validation.png
% =========================================================================

\section{Integrated Multi-Mechanism System Fragility: Historical Climate Baseline}
\label{sec: Integrated Multi-Mechanism System Fragility: Historical Climate Baseline}

The series-system composition of Section~\ref{sec: Multi-Mechanism
Integration: Series-System Joint-Probability Formulation} was executed for
all 114 segments of the study reaches: the posterior transient BEP
fragility of Phase~2 at the four borehole-anchored sections, and the
overflow and fluvial-scour fragility curves at every segment. The surface
curves were obtained by faithful re-execution of the failure-judgment
models of \textcite{uemura_phd_2025} and \textcite{uemura_wp2_2024} on
Uemura's own consolidated per-segment input table, whose stage--discharge
coefficients are numerically identical to the rating relations of
Section~\ref{sec: The d4PDF Climate Ensemble: Structure, Downscaling, and
Hydrograph Extraction}, so that all three mechanisms share one hydraulic
axis by construction. Each curve conditions on the same canonical d4PDF
event shape that conditioned the BEP fragility, scaled per stage level,
with $10^4$ Monte Carlo realizations of the published uncertainty models
(rating error, crest-height error, and turf critical velocity for
overflow, per \textcite{dean_2010}; erodibility and critical shear stress
for scour) under common random numbers, which renders every conditional
curve exactly non-decreasing. The re-execution route, its seeds, and its
validation gates are recorded in the repository decision log (ADR-0042);
two findings of that re-execution --- a unit-conversion discrepancy in the
scour erodibility coefficient whose dimensionally consistent correction
suppresses the scour mechanism entirely, and a regularized singularity of
the friction-factor log law at millimetre flow depths --- are carried as
flagged companions rather than silent repairs, and the primary curves
reproduce the source implementation as received.

Figure~\ref{fig:phase3-bep-sections} shows the composed conditional system
fragility $P_{f,\mathrm{segment}}(h)$ at the four BEP sections under the
historical baseline. The structure is consistent across KP~57.4, KP~58.8,
and KP~60.0: the posterior BEP transition rises one to two metres below
the stage at which overflow becomes possible (crest exceedance within the
rating-error scatter), so the subsurface mechanism controls the lower
shoulder of the system curve, and overflow joins only near and above the
crest. KP~62.0 differs in degree rather than in
order: its transient BEP transition is bracketed only within the
hypothetical above-crest extension of its conditioning grid, so that part
of the curve must never be read as attainable, but the subsurface
mechanism still controls the system curve there, carrying 81~per cent of
the summed annual contribution historically and half of it under $+4$~K.
Annualized against the historical (HPB, 3{,}000-year) empirical
peak-stage distribution of each node, the segment failure probabilities at
the four sections are $7.5\times10^{-4}$, $7.4\times10^{-3}$,
$1.8\times10^{-3}$, and $1.0\times10^{-3}$~yr$^{-1}$ respectively, and the
basin-wide maximum over all 114 segments is the BEP-governed segment at
KP~58.8.

\begin{figure}[htbp]
    \centering
    \includegraphics[width=\textwidth]{figures/phase3_system_fragility_bep_sections.png}
    \caption[Composed three-mechanism segment fragility at the BEP
    sections]{Composed three-mechanism conditional segment fragility at
    the four borehole-anchored sections (posterior BEP, matrix $d_{70}$,
    $\lambda_{ac}=250$~m). Shaded bands mark stages beyond the largest
    +4K ensemble peak; the KP~62.0 grid extension is a fit stabilizer,
    not attainable loading.}
    \label{fig:phase3-bep-sections}
\end{figure}

\section{Relative Dominance of Failure Mechanisms under Historical Loading}
\label{sec: Relative Dominance of Failure Mechanisms under Historical Loading}

Table~\ref{tab:phase3-dominance-hist} gives the annualized per-mechanism
contributions and dominance shares at the BEP sections. Backward erosion
piping is the dominant mechanism at three of the four sections where it is
quantified, carrying 88--99.9\% of the summed annual contributions at
KP~57.4, KP~58.8, and KP~60.0; overflow dominates only at KP~62.0 (79\%).
Fluvial scour never exceeds 9\% --- a share computed under the source
implementation's fast erodibility conversion, and identically zero under
the dimensionally consistent correction, so its true contribution is
bounded by a value that is already immaterial to the system.

\begin{table}[htbp]
    \centering
    \caption[Annualized mechanism dominance at the BEP sections,
    historical baseline]{Annualized per-mechanism failure probabilities
    [yr$^{-1}$] and dominance shares under the historical baseline
    (posterior BEP, matrix $d_{70}$).}
    \label{tab:phase3-dominance-hist}
    \begin{tabular}{lcccc}
        \toprule
        Section & $P_{f,\mathrm{sys}}$ & BEP & Overflow & Fluvial scour \\
        \midrule
        KP~57.4 & $8.1\times10^{-4}$ & $7.5\times10^{-4}$ (91\%) & $3\times10^{-6}$ (0.4\%) & $7\times10^{-5}$ (8\%) \\
        KP~58.8 & $7.9\times10^{-3}$ & $7.3\times10^{-3}$ (88\%) & $8.3\times10^{-4}$ (10\%) & $2\times10^{-4}$ (3\%) \\
        KP~60.0 & $1.8\times10^{-3}$ & $1.8\times10^{-3}$ (99.9\%) & 0 (exact) & $2\times10^{-6}$ (0.1\%) \\
        KP~62.0 & $2.6\times10^{-3}$ & $3.5\times10^{-4}$ (12\%) & $2.2\times10^{-3}$ (79\%) & $2.6\times10^{-4}$ (9\%) \\
        \bottomrule
    \end{tabular}
\end{table}

Along the full reaches (Figure~\ref{fig:phase3-dominance-profile}), the
110 segments without borehole anchoring compose the two surface mechanisms
only, and their system values are therefore lower bounds that omit the
mechanism found dominant wherever it could be quantified. This gap is
material rather than cosmetic --- the basin's single riskiest segment is
BEP-governed --- and motivates the bounded-extrapolation prior extension
for the borehole-free reaches identified as future work in
Chapter~\ref{chap: Discussion}. Aggregated to Uemura's consequence
sections under his full-dependence rule (the within-section maximum taken
conditional on discharge, mapped through each member segment's own rating
relation), section Tokachi~4 (KP~58.0) governs the basin at
$8.0\times10^{-3}$~yr$^{-1}$, and its governing member is the BEP-dominant
segment KP~58.8: the subsurface mechanism controls the worst consequence
section of the system.

\begin{figure}[htbp]
    \centering
    \includegraphics[width=\textwidth]{figures/phase3_dominance_profile.png}
    \caption[Longitudinal annualized mechanism profile]{Annualized
    per-mechanism failure probability along both study reaches and both
    scenarios (posterior BEP, matrix $d_{70}$). BEP exists at the four
    isolated borehole sections (markers); values at the display floor are
    exact zeros.}
    \label{fig:phase3-dominance-profile}
\end{figure}

\section{Integrated Multi-Mechanism System Fragility: \texorpdfstring{+4K}{+4K} Warming Scenario}
\label{sec: Integrated Multi-Mechanism System Fragility: +4K Warming Scenario}

Under the +4K hazard (HFB, 5{,}400 years) the identical conditional
fragility surfaces --- climate enters only through the stage-frequency, per
the shape-invariance established for the BEP curves and imposed
structurally for the surface curves --- produce the annualized values of
Figure~\ref{fig:phase3-climate-shift}. The median annual system failure
probability across the 114 segments rises from $6.3\times10^{-5}$ to
$2.8\times10^{-3}$~yr$^{-1}$, a factor 44; the number of segments above
$10^{-3}$~yr$^{-1}$ grows from 14 to 80, and above $10^{-2}$~yr$^{-1}$
from 0 to 19. At the BEP sections the system ratios are 5.4--12.6, and the
basin maximum remains the BEP-governed KP~58.8 at
$4.3\times10^{-2}$~yr$^{-1}$. The sensitivity brackets carried alongside
the primary numbers behave as the mechanism structure predicts: the
$\lambda_{ac}$ length-effect bracket (40~m floor,
\textcite{kanning_2012}) multiplies the BEP-driven system numbers by up to
3.4 but leaves overflow-dominated KP~62.0 nearly unchanged; the resistant
bulk-$d_{70}$ reading cuts the BEP-driven numbers 4--25$\times$; and the
2016 survival constraint of Phase~2 lowers the KP~58.8 historical system
value by 11\% relative to the prior.

\begin{figure}[htbp]
    \centering
    \includegraphics[width=\textwidth]{figures/phase3_climate_shift.png}
    \caption[Climate shift of the annualized system failure
    probability]{Annualized system failure probability per segment,
    historical versus +4K, with the +4K/historical ratio (lower panels)
    and the $\lambda_{ac}=40$~m bracket at the BEP sections.}
    \label{fig:phase3-climate-shift}
\end{figure}

\section{Climate-Induced Shift in Failure Mode Dominance}
\label{sec: Climate-Induced Shift in Failure Mode Dominance}

The warming scenario does not displace the subsurface mechanism where it
is quantified: BEP retains 74--95\% dominance at KP~57.4, KP~58.8, and
KP~60.0 under +4K. What shifts is the surface picture around it. First,
overflow's share grows wherever crests come within reach of the amplified
peaks --- most sharply at KP~57.4, from 0.4\% to 20\% of the summed
contributions. Second, among the surface-only segments the reach-wide
dominant surface mechanism flips from fluvial scour (71 of 114 segments
historically) to overflow (82 of 114 under +4K), because crest exceedance
turns from a tail event into a recurrent one. Third, and most relevant to
the adaptation question, the attribution of
Figure~\ref{fig:phase3-rq4-attribution} shows the climate signal traveling
through the duration channel: annual system risk concentrates roughly two
orders of magnitude in years that load the landside toe for more than
24~hours (e.g.\ KP~58.8 historical: conditional $P_f$ 0.138 in the
long-loading stratum versus $9.7\times10^{-4}$ outside it), and the +4K
ensemble roughly triples the frequency of exactly those years (5.1\% to
13.5\% of years at KP~58.8). Compound years with two or more separate
above-toe excursions sit between the two duration strata at every
section. The mechanism-dominance answer to the research question is
therefore asymmetric: climate change amplifies every mechanism's absolute
risk, promotes overflow from negligible to a first-order contributor at
the near-crest sections, and yet leaves backward erosion piping the
governing mechanism of the basin's worst segments --- because the same
elongated-duration years that enable overflow are the years in which the
transient BEP mechanism has time to progress to breach.

\begin{figure}[htbp]
    \centering
    \includegraphics[width=\textwidth]{figures/phase3_rq4_attribution.png}
    \caption[Duration- and compound-stratified conditional failure
    probability]{RQ4 attribution at the BEP sections: conditional annual
    system failure probability within duration and compound-event strata,
    historical versus +4K (posterior BEP, matrix $d_{70}$).}
    \label{fig:phase3-rq4-attribution}
\end{figure}

The composition inherits, and states, the limitations of its inputs: the
conditional-independence-given-$h$ assumption of \textcite{pol_ress_2023}
excludes dynamic mechanism coupling (non-conservative, as discussed in
Chapter~\ref{chap: Discussion}); the drained sections are evaluated as
unremediated foundations; KP~62.0's BEP share is a grid-clamped lower
bound; and the canonical-shape conditioning shared by all three mechanisms
was quantified against a full event-based re-execution of the surface
models over both ensembles at the nine section-representative nodes ---
essentially exact for overflow (agreement within tens of percent) and
2--15$\times$ conservative for the duration-integrating scour mechanism,
whose contribution to the system is in any case bounded by the
unit-conversion finding above.
